\subsection{TOPOLOGICAL RINGS}

DEFINITION 2. A topological ring is a set A which carries a ring structure and a topology satisfying the following axioms:

\((\mathrm{AT}_{\mathbf{I}})\) . The mapping \((x,y) \to x + y\) of \(\mathbf{A} \times \mathbf{A}\) into \(\mathbf{A}\) is continuous.

\((\mathrm{AT}_{\mathbf{II}})\) . The mapping \(x \to -x\) of A into A is continuous.

\((\mathrm{AT}_{\mathrm{III}})\) . The mapping \((x,y)\to xy\) of \(\mathbf{A}\times \mathbf{A}\) into \(\mathbf{A}\) is continuous.

The first two axioms express that the topology of \(\mathbf{A}\) is compatible with its additive group structure (§ 1, no. 1).

If a ring structure and a topology are given on a set A, they are said to be compatible if they satisfy axioms \((\mathrm{AT}_{\mathbf{I}})\) , \((\mathrm{AT}_{\mathbf{II}})\) and \((\mathrm{AT}_{\mathbf{III}})\) .

Examples. 1) On any ring A, the discrete topology is compatible with the ring structure. A topological ring whose topology is discrete is called a discrete ring.

* 2) We shall see in Chapter IV that the topology of the rational line Q (resp. the real line R) is compatible with the ring structure of Q (resp. R). *

In a topological ring every left homothety \(x \rightarrow ax\) (resp. every right homothety \(x \rightarrow xa\) ) is continuous (and is a homeomorphism if a is a unit of A).

Since we can write identically

\[
x y - x _ {0} y _ {0} = (x - x _ {0}) (y - y _ {0}) + (x - x _ {0}) y _ {0} + x _ {0} (y - y _ {0})
\]

the axiom \((\mathrm{AT}_{\mathrm{III}})\) {[}in view of \((\mathrm{AT}_{\mathrm{I}})\) and \((\mathrm{AT}_{\mathrm{II}})\) {]} is equivalent to the conjunction of the following two axioms:

\((\mathrm{AT}_{\mathrm{III}a})\) . Given any \(x_0\in \mathbf{A}\) , the mappings \(x\to x_0x\) and \(x\rightarrow xx_0\) are continuous at the point \(x = 0\)

\((\mathrm{AT}_{\mathbf{III}\ b})\) . The mapping \((x,y)\rightarrow xy\) of \(\mathbf{A}\times \mathbf{A}\) into A is continuous at the point (o, o).

From this we can deduce a necessary and sufficient set of conditions which the filter \(\mathfrak{B}\) of neighbourhoods of \(o\) in a ring \(A\) must satisfy in order to define a topology on \(A\) compatible with its ring structure: \(\mathfrak{B}\) must satisfy axioms \((\mathrm{GA}_{\mathrm{I}})\) and \((\mathrm{GA}_{\mathrm{II}})\) of § 1, and also the following two axioms:

\((\mathrm{AV}_{1})\) . For all \(x_0\in \mathbf{A}\) and all \(\mathbf{V}\in \mathfrak{B}\) , there exists \(\mathbf{W}\in \mathfrak{B}\) such that \(x_0\mathbf{W}\subset \mathbf{V}\) and \(\mathbf{W}x_0\subset \mathbf{V}\)

\((\mathrm{AV}_{\mathrm{II}})\) . For all \(\mathbf{V} \in \mathfrak{B}\) , there exists \(\mathbf{W} \in \mathfrak{B}\) such that \(\mathbf{WW} \subset \mathbf{V}\) .

Remark. In analysis one fairly often meets rings which satisfy axioms \((\mathrm{AT}_{\mathrm{I}})\) , \((\mathrm{AT}_{\mathrm{II}})\) and \((\mathrm{AT}_{\mathrm{III}a})\) , but not \((\mathrm{AT}_{\mathrm{III}b})\) . * An example is the ring of measures on a compact group, where multiplication is convolution and the topology is the vague topology. *

Example 3). Let \(\mathfrak{B}\) be a filter base on a ring A, consisting of two-sided ideals. \(\mathfrak{B}\) is a fundamental system of neighbourhoods of o for a topology compatible with the additive group structure of A, and it follows immediately from \((\mathrm{AV}_{\mathrm{I}})\) and \((\mathrm{AV}_{\mathrm{II}})\) that this topology is compatible with the ring structure of A.

Let X be a topological space, and let f and g be two mappings of X into a topological ring A. If f and g are continuous at a point \(x_{0} \in X\) , then \(f + g\) , -f and fg are continuous at this point. It follows that the continuous mappings of X into A form a subring of the ring \(A^{X}\) of all mappings of X into A. We see also that, if A is commutative, then every polynomial in n variables, with coefficients in A and defined on \(A^{n}\) , is continuous on \(A^{n}\) . Again, let f and g be two mappings of a set X, filtered by a filter F, into a Hausdorff topological ring A; if \(\lim_{\mathfrak{F}} f\) and \(\lim_{\mathfrak{F}} g\) exist, then so do \(\lim_{\mathfrak{F}} (f + g)\) , \(\lim_{\mathfrak{F}} (-f)\) and \(\lim_{\mathfrak{F}} (fg)\) , and we have (Chapter I, §7, no. 4, Proposition 9, Corollary 1, and §8, no. 1, Proposition 1)

\[
\lim _ {\mathfrak {F}} (f + g) = \lim _ {\mathfrak {F}} f + \lim _ {\mathfrak {F}} g,\tag{1}
\]

\[
\lim _ {\mathfrak {F}} (- f) = - \lim _ {\mathfrak {F}} f,\tag{2}
\]

\[
\lim _ {\mathfrak {F}} (f g) = (\lim _ {\mathfrak {F}} f) (\lim _ {\mathfrak {F}} g).\tag{3}
\]

